\documentclass[12pt]{article}
\usepackage[ukrainian]{babel}
\usepackage{fontspec}
% --- НАЛАШТУВАННЯ ШРИФТІВ ---
\setmainfont{Schoolbook-Regular.otf}[
    Path = ../../,
    BoldFont = Schoolbook-Bold.otf,
    ItalicFont = Schoolbook-Italic.otf,
    BoldItalicFont = Schoolbook-BoldItalic.otf
]
\usepackage[italic]{mathastext}
\usepackage[a4paper,margin=1.8cm,bottom=2cm,top=2cm]{geometry}
\usepackage{amsmath,amssymb}
\usepackage{tikz}
\usepackage{pgfplots}
\pgfplotsset{compat=1.16}
\usetikzlibrary{calc,patterns,angles,quotes}
\usepackage{xcolor,array,fancyhdr,enumitem,multicol}
\usepackage{colortbl}
\usepackage{tcolorbox}
\tcbuselibrary{skins}
\usepackage[hidelinks]{hyperref}
\usepackage[firstpage=false, text=@pvtr2525, scale=0.6, color=gray!12]{draftwatermark}
\definecolor{mainGreen}{RGB}{34, 120, 64}
\definecolor{yearOrange}{RGB}{220, 100, 30}
\definecolor{yearcolor}{RGB}{220, 100, 30}
\definecolor{titleBg}{RGB}{225, 240, 220}
\definecolor{instrBg}{RGB}{255, 240, 220}
\definecolor{instrBorder}{RGB}{220, 150, 80}
\pagestyle{fancy}
\fancyhf{}
\renewcommand{\headrulewidth}{0.4pt}
\renewcommand{\footrulewidth}{0.4pt}
\fancyhead[C]{\small\color{gray!80} Збірник НМТ \textendash{} сесії 23.05--5.06.2026}
\fancyhead[R]{\small\color{mainGreen}\textbf{\href{https://t.me/pvtr2525}{@pvtr2525}}}
\fancyfoot[L]{\small\color{gray!80}\href{https://t.me/pvtr2525}{ТГ: @pvtr2525}}
\fancyfoot[C]{\small\color{gray!80}\thepage}
\fancyfoot[R]{\small\color{gray!80}НМТ 2026}
% --- ТАБЛИЦІ ВІДПОВІДЕЙ ---
\newcommand{\answerTable}[5]{%
\par\vspace{0.15cm}
\noindent
\begin{tabular}{|*{5}{>{\centering\arraybackslash}m{3cm}|}}
\hline
\rule[-0.15cm]{0pt}{0.6cm}\textbf{А} & \rule[-0.15cm]{0pt}{0.6cm}\textbf{Б} & \rule[-0.15cm]{0pt}{0.6cm}\textbf{В} & \rule[-0.15cm]{0pt}{0.6cm}\textbf{Г} & \rule[-0.15cm]{0pt}{0.6cm}\textbf{Д} \\
\hline
\rule[-0.2cm]{0pt}{0.75cm}#1 & \rule[-0.2cm]{0pt}{0.75cm}#2 & \rule[-0.2cm]{0pt}{0.75cm}#3 & \rule[-0.2cm]{0pt}{0.75cm}#4 & \rule[-0.2cm]{0pt}{0.75cm}#5 \\
\hline
\end{tabular}
\par\vspace{0.25cm}
}
\newcommand{\answerTableTall}[5]{%
\par\vspace{0.15cm}
\noindent
\begin{tabular}{|*{5}{>{\centering\arraybackslash}m{3cm}|}}
\hline
\rule[-0.2cm]{0pt}{0.7cm}\textbf{А} & \rule[-0.2cm]{0pt}{0.7cm}\textbf{Б} & \rule[-0.2cm]{0pt}{0.7cm}\textbf{В} & \rule[-0.2cm]{0pt}{0.7cm}\textbf{Г} & \rule[-0.2cm]{0pt}{0.7cm}\textbf{Д} \\
\hline
\rule[-0.45cm]{0pt}{1.2cm}#1 & \rule[-0.45cm]{0pt}{1.2cm}#2 & \rule[-0.45cm]{0pt}{1.2cm}#3 & \rule[-0.45cm]{0pt}{1.2cm}#4 & \rule[-0.45cm]{0pt}{1.2cm}#5 \\
\hline
\end{tabular}
\par\vspace{0.25cm}
}
\newcommand{\instructionBox}[1]{%
\par\vspace{0.3cm}
\begin{tcolorbox}[colback=instrBg, colframe=instrBorder, boxrule=0.6pt, arc=2pt,
    left=8pt, right=8pt, top=4pt, bottom=4pt]
\centering\small\bfseries #1
\end{tcolorbox}
\par\vspace{0.15cm}
}
\newcommand{\sectionTitle}[1]{%
\par\vspace{0.4cm}
\begin{tcolorbox}[colback=titleBg, colframe=titleBg, boxrule=0pt, arc=8pt,
    halign=center, fontupper=\Large\bfseries\color{mainGreen}]
#1
\end{tcolorbox}
\par\vspace{0.3cm}
}
\newcommand{\task}[2]{\noindent\textbf{#1.}\ \ #2\par\vspace{0.2cm}}
% --- НМТ-СТИЛЬ ПОЛЕ ВІДПОВІДІ ---
\newcommand{\nmtAnswerBox}{%
\par\vspace{0.2cm}\noindent
Відповідь:\ \
\begin{tabular}{@{}*{4}{|>{\centering\arraybackslash}p{0.8cm}}|@{\hspace{0.1cm}\textbf{,}\hspace{0.1cm}}*{3}{|>{\centering\arraybackslash}p{0.8cm}}|@{}}
\hline
\rule[-0.2cm]{0pt}{0.7cm} & & & & & & \\
\hline
\end{tabular}
\par\vspace{0.3cm}
}
% --- СІТКА ДЛЯ МАТЧИНГУ ---
\newcommand{\matchingGrid}{%
    \begingroup
    \renewcommand{\arraystretch}{1.15}
    \setlength{\tabcolsep}{4pt}
    \footnotesize
    \begin{tabular}{r|c|c|c|c|c|}
         \multicolumn{1}{c}{} & \multicolumn{1}{c}{\textbf{А}} & \multicolumn{1}{c}{\textbf{Б}} & \multicolumn{1}{c}{\textbf{В}} & \multicolumn{1}{c}{\textbf{Г}} & \multicolumn{1}{c}{\textbf{Д}} \\ \cline{2-6}
         \textbf{1} & \phantom{X} & \phantom{X} & \phantom{X} & \phantom{X} & \phantom{X} \\ \cline{2-6}
         \textbf{2} & & & & & \\ \cline{2-6}
         \textbf{3} & & & & & \\ \cline{2-6}
    \end{tabular}
    \endgroup
}
% --- ДОДАТКОВІ МАКРОСИ ЗБІРНИКА ---
\newcommand{\taskBlock}[1]{%
\par\vspace{0.45cm}
\noindent{\color{mainGreen}\rule{\linewidth}{0.8pt}}\par\vspace{0.12cm}
\noindent{\small\bfseries\color{mainGreen}#1}\par\vspace{0.25cm}
}
\newcounter{zad}
\newcommand{\zadtask}[1]{\stepcounter{zad}\noindent\textbf{\thezad.}\ \ #1\par\vspace{0.2cm}}
\newcommand{\zadnum}{\stepcounter{zad}\textbf{\thezad.}\ \ }
\newcounter{ans}
\newcommand{\ansitem}[1]{\stepcounter{ans}\noindent\textbf{\theans.}\ #1\par\vspace{0.07cm}}
\newcommand{\nmtyear}[1]{\hfill{\small\color{yearOrange}(НМТ #1)}}% --- ДОДАТКОВІ МАКРОСИ (розв'язки, зміст) ---
\tcbuselibrary{breakable}
\newcommand{\solution}[1]{%
\par\vspace{0.12cm}
\begin{tcolorbox}[colback=titleBg!45, colframe=mainGreen!55, boxrule=0.4pt, arc=2pt,
    left=6pt, right=6pt, top=3pt, bottom=3pt, breakable]
\footnotesize\textbf{Короткий розв'язок (оригінал).}\ #1
\end{tcolorbox}
\par\vspace{0.12cm}
}
\setcounter{tocdepth}{2}
\newcommand{\chapterTitle}[1]{%
\clearpage\phantomsection\addcontentsline{toc}{section}{#1}%
\sectionTitle{#1}}
\newcommand{\typeTitle}[1]{%
\phantomsection\addcontentsline{toc}{subsection}{\quad #1}%
\par\vspace{0.3cm}
\noindent{\large\bfseries\color{yearOrange}#1}\par\vspace{0.08cm}
\noindent{\color{yearOrange}\rule{\linewidth}{0.4pt}}\par\vspace{0.2cm}}
\newcommand{\ansTheme}[1]{\par\vspace{0.25cm}\noindent{\bfseries\color{mainGreen}#1}\par\vspace{0.12cm}}
\newcommand{\ansType}[1]{\par\vspace{0.1cm}\noindent{\itshape\color{yearOrange}#1}\par\vspace{0.08cm}}
\begin{document}
\thispagestyle{empty}
\vspace*{1.2cm}
\begin{center}
{\Large\bfseries\color{mainGreen} ЗБІРНИК НМТ-2026 з математики}\\[0.4cm]
{\Huge\bfseries\color{mainGreen} РОЗДІЛ 3. ФУНКЦІЇ}\\[0.6cm]
{\large Оригінали сесій 23.05--5.06.2026 з короткими розв'язками та авторськими аналогами}\\[1cm]
{\large\color{mainGreen}\textbf{\href{https://t.me/pvtr2525}{Телеграм: @pvtr2525}}}
\end{center}
\clearpage
\phantomsection\addcontentsline{toc}{section}{Зміст}
\sectionTitle{ЗМІСТ}
{\hypersetup{linkcolor=black}\renewcommand{\contentsname}{}\vspace{-1cm}\tableofcontents}
\chapterTitle{РОЗДІЛ 3. ФУНКЦІЇ}
\typeTitle{Завдання з вибором однієї відповіді (А--Д)}
\taskBlock{Зсув графіка функції, координатні чверті --- сесія 23.05, завдання №4}
\noindent
\begin{minipage}[c]{0.55\textwidth}
\zadnum На рисунку зображено графік функції $y = f(x)$. У яких координатних чвертях розміщений графік функції $y = f(x) + 3$? \nmtyear{2026}
\end{minipage}%
\hfill
\begin{minipage}[c]{0.43\textwidth}
\centering
\begin{tikzpicture}[scale=0.55, thick]
    \draw[->] (-4,0) -- (4.5,0) node[right] {$x$};
    \draw[->] (0,-3) -- (0,5) node[above] {$y$};
    \node[below] at (-3,-0.1) {\scriptsize $-3$};
    \node[below] at (1,-0.1) {\scriptsize $1$};
    \node[below] at (3,-0.1) {\scriptsize $3$};
    \node[left] at (-0.1,4) {\scriptsize $4$};
    \node[left] at (-0.1,-2) {\scriptsize $-2$};
    \draw[domain=-3:3, smooth, variable=\x, thick]
        plot ({\x},{4 - 0.49*(\x - 0.5)*(\x - 0.5)});
    \fill (-3,-2) circle (2pt);
    \fill (3,0.94) circle (2pt);
    \draw[dashed] (-3,-2) -- (-3,0);
    \draw[dashed] (-3,-2) -- (0,-2);
    \draw[dashed] (3,0.94) -- (3,0);
    \node[above right] at (1.6,2.5) {\scriptsize $f(x)$};
\end{tikzpicture}
\end{minipage}
\par\vspace{0.3cm}
\noindent
\textbf{А)} \ лише у I та II \\[0.15cm]
\textbf{Б)} \ лише у III та IV \\[0.15cm]
\textbf{В)} \ лише у I та III \\[0.15cm]
\textbf{Г)} \ лише у II та IV \\[0.15cm]
\textbf{Д)} \ в усіх чотирьох
\par\vspace{0.2cm}

\solution{Графік $y=f(x)$ має значення від $-2$ до $4$. Зсув на $+3$ угору дає значення від $1$ до $7$~--- усі додатні. Оскільки $x$ набуває як від'ємних, так і додатних значень, графік лежить лише в I та II чвертях. Відповідь: \textbf{А}.}

\noindent
\begin{minipage}[c]{0.55\textwidth}
\zadnum На рисунку зображено графік функції $y = f(x)$. У яких координатних чвертях розміщений графік функції $y = f(x) - 6$?
\end{minipage}%
\hfill
\begin{minipage}[c]{0.43\textwidth}
\centering
\begin{tikzpicture}[scale=0.55, thick]
    \draw[->] (-4,0) -- (4.5,0) node[right] {$x$};
    \draw[->] (0,-1.5) -- (0,6) node[above] {$y$};
    \node[below] at (-3,-0.1) {\scriptsize $-3$};
    \node[below] at (3,-0.1) {\scriptsize $3$};
    \node[left] at (-0.1,5) {\scriptsize $5$};
    \node[left] at (-0.1,1) {\scriptsize $1$};
    \draw[domain=-3:3, smooth, variable=\x, thick]
        plot ({\x},{5 - 0.444*\x*\x});
    \fill (-3,1) circle (2pt);
    \fill (3,1) circle (2pt);
    \draw[dashed] (-3,1) -- (-3,0);
    \draw[dashed] (3,1) -- (3,0);
    \draw[dashed] (-3,1) -- (0,1);
    \node[above right] at (1.4,3.2) {\scriptsize $f(x)$};
\end{tikzpicture}
\end{minipage}
\par\vspace{0.3cm}
\noindent
\textbf{А)} \ лише у I та II \\[0.15cm]
\textbf{Б)} \ лише у III та IV \\[0.15cm]
\textbf{В)} \ лише у I та III \\[0.15cm]
\textbf{Г)} \ лише у II та IV \\[0.15cm]
\textbf{Д)} \ в усіх чотирьох
\par\vspace{0.2cm}

\noindent
\begin{minipage}[c]{0.55\textwidth}
\zadnum На рисунку зображено графік функції $y = f(x)$. У яких координатних чвертях розміщений графік функції $y = f(x) + 6$?
\end{minipage}%
\hfill
\begin{minipage}[c]{0.43\textwidth}
\centering
\begin{tikzpicture}[scale=0.55, thick]
    \draw[->] (-4,0) -- (4.5,0) node[right] {$x$};
    \draw[->] (0,-5) -- (0,1.5) node[above] {$y$};
    \node[below] at (-3,-0.1) {\scriptsize $-3$};
    \node[below] at (3,-0.1) {\scriptsize $3$};
    \node[left] at (-0.1,-4) {\scriptsize $-4$};
    \node[left] at (-0.1,-1) {\scriptsize $-1$};
    \draw[domain=-3:3, smooth, variable=\x, thick]
        plot ({\x},{-4 + 0.333*\x*\x});
    \fill (-3,-1) circle (2pt);
    \fill (3,-1) circle (2pt);
    \draw[dashed] (-3,-1) -- (-3,0);
    \draw[dashed] (3,-1) -- (3,0);
    \draw[dashed] (-3,-1) -- (0,-1);
    \node[below right] at (1.2,-2.6) {\scriptsize $f(x)$};
\end{tikzpicture}
\end{minipage}
\par\vspace{0.3cm}
\noindent
\textbf{А)} \ лише у I та II \\[0.15cm]
\textbf{Б)} \ лише у III та IV \\[0.15cm]
\textbf{В)} \ лише у I та III \\[0.15cm]
\textbf{Г)} \ лише у II та IV \\[0.15cm]
\textbf{Д)} \ в усіх чотирьох
\par\vspace{0.2cm}
\taskBlock{Член послідовності за формулою --- сесія 23.05, завдання №6}
\zadtask{Послідовність задано формулою $a_n = (-2)^n - 7$. Визначте її четвертий член. \nmtyear{2026}}
\answerTable{$-23$}{$23$}{$9$}{$-15$}{$1$}
\solution{Підставимо $n=4$: $a_4=(-2)^4-7=16-7=9$. Відповідь: \textbf{В}.}
\zadtask{Послідовність задано формулою $a_n = (-3)^n + 5$. Визначте її третій член.}
\answerTable{$32$}{$-22$}{$22$}{$-4$}{$14$}
\zadtask{Послідовність задано формулою $a_n = \dfrac{n^2 - 1}{n + 3}$. Визначте її п'ятий член.}
\answerTable{$3$}{$\dfrac{25}{8}$}{$4$}{$\dfrac{8}{3}$}{$2$}
\taskBlock{Паралельність графіків лінійних функцій --- сесія 30.05, завдання №8}
\noindent
\begin{minipage}[c]{0.55\textwidth}
\zadnum На рисунку зображено графік функції $y = -2x$. Графік якої з наведених функцій паралельний графіку цієї функції?
\end{minipage}%
\hfill
\begin{minipage}[c]{0.43\textwidth}
\centering
\begin{tikzpicture}[scale=0.45, thick]
    \draw[gray!30, very thin] (-4.5,-4.5) grid (4.5,4.5);
    \draw[->, thick] (-4.7,0) -- (4.7,0) node[right] {\small $x$};
    \draw[->, thick] (0,-4.7) -- (0,4.7) node[above] {\small $y$};
    \foreach \x in {-4,-3,-2,-1,1,2,3,4} {
        \draw[thick] (\x,-0.1) -- (\x,0.1);
        \node[below] at (\x,-0.1) {\scriptsize $\x$};
    }
    \foreach \y in {-4,-3,-2,-1,1,2,3,4} {
        \draw[thick] (-0.1,\y) -- (0.1,\y);
        \node[left] at (-0.1,\y) {\scriptsize $\y$};
    }
    \node[below left] at (0,0) {\scriptsize $0$};
    \draw[thick, mainGreen, line width=1.2pt] (-2.2,4.4) -- (2.2,-4.4);
    \node[mainGreen, fill=white, inner sep=1pt] at (2.6,-3.5) {\scriptsize $y = -2x$};
\end{tikzpicture}
\end{minipage}
\par\vspace{0.2cm}
\answerTable{$y = 2x + 1$}{$y = 1 - 2x$}{$y = -x - 2$}{$y = \dfrac{x}{2} - 1$}{$y = -\dfrac{1}{2}x$}
\solution{Прямі паралельні, коли кутові коефіцієнти рівні. У $y=-2x$ маємо $k=-2$, тож шукаємо функцію з $k=-2$: це $y=1-2x$. Відповідь: \textbf{Б}.}
\noindent
\begin{minipage}[c]{0.55\textwidth}
\zadnum На рисунку зображено графік функції $y = 3x$. Графік якої з наведених функцій паралельний графіку цієї функції?
\end{minipage}%
\hfill
\begin{minipage}[c]{0.43\textwidth}
\centering
\begin{tikzpicture}[scale=0.45, thick]
    \draw[gray!30, very thin] (-4.5,-4.5) grid (4.5,4.5);
    \draw[->, thick] (-4.7,0) -- (4.7,0) node[right] {\small $x$};
    \draw[->, thick] (0,-4.7) -- (0,4.7) node[above] {\small $y$};
    \foreach \x in {-4,-3,-2,-1,1,2,3,4} {
        \draw[thick] (\x,-0.1) -- (\x,0.1);
        \node[below] at (\x,-0.1) {\scriptsize $\x$};
    }
    \foreach \y in {-4,-3,-2,-1,1,2,3,4} {
        \draw[thick] (-0.1,\y) -- (0.1,\y);
        \node[left] at (-0.1,\y) {\scriptsize $\y$};
    }
    \node[below left] at (0,0) {\scriptsize $0$};
    \draw[thick, mainGreen, line width=1.2pt] (-1.45,-4.35) -- (1.45,4.35);
    \node[mainGreen, fill=white, inner sep=1pt] at (2.6,3.5) {\scriptsize $y = 3x$};
\end{tikzpicture}
\end{minipage}
\par\vspace{0.2cm}
\answerTable{$y = -3x$}{$y = \dfrac{x}{3} + 2$}{$y = 3x - 2$}{$y = -\dfrac{x}{3}$}{$y = x + 3$}
\noindent
\begin{minipage}[c]{0.55\textwidth}
\zadnum На рисунку зображено графік функції $y = \dfrac{x}{2}$. Графік якої з наведених функцій \textit{перпендикулярний} до графіка цієї функції?
\end{minipage}%
\hfill
\begin{minipage}[c]{0.43\textwidth}
\centering
\begin{tikzpicture}[scale=0.45, thick]
    \draw[gray!30, very thin] (-4.5,-4.5) grid (4.5,4.5);
    \draw[->, thick] (-4.7,0) -- (4.7,0) node[right] {\small $x$};
    \draw[->, thick] (0,-4.7) -- (0,4.7) node[above] {\small $y$};
    \foreach \x in {-4,-3,-2,-1,1,2,3,4} {
        \draw[thick] (\x,-0.1) -- (\x,0.1);
        \node[below] at (\x,-0.1) {\scriptsize $\x$};
    }
    \foreach \y in {-4,-3,-2,-1,1,2,3,4} {
        \draw[thick] (-0.1,\y) -- (0.1,\y);
        \node[left] at (-0.1,\y) {\scriptsize $\y$};
    }
    \node[below left] at (0,0) {\scriptsize $0$};
    \draw[thick, mainGreen, line width=1.2pt] (-4.4,-2.2) -- (4.4,2.2);
    \node[mainGreen, fill=white, inner sep=1pt] at (3.3,2.9) {\scriptsize $y = \frac{x}{2}$};
\end{tikzpicture}
\end{minipage}
\par\vspace{0.2cm}
\answerTable{$y = 2x + 1$}{$y = -\dfrac{x}{2}$}{$y = 2 - 2x$}{$y = \dfrac{x}{2} + 3$}{$y = -2x - 1$}
\taskBlock{Порівняння інтегралів --- сесія 30.05, завдання №11}
\zadtask{Укажіть інтеграл, який має \textit{найменше} значення.

\par\vspace{0.2cm}
\begin{flushleft}
\textbf{А} \quad $\displaystyle\int_0^1 \dfrac{1}{4}\,dx$ \\[0.25cm]
\textbf{Б} \quad $\displaystyle\int_0^1 \dfrac{1}{3}\,dx$ \\[0.25cm]
\textbf{В} \quad $\displaystyle\int_0^1 \dfrac{1}{2}\,dx$ \\[0.25cm]
\textbf{Г} \quad $\displaystyle\int_0^1 1\,dx$ \\[0.25cm]
\textbf{Д} \quad $\displaystyle\int_0^1 \sqrt{2}\,dx$
\end{flushleft}
}
\solution{Для сталої $c$ маємо $\int_0^1 c\,dx=c$. Тож значення дорівнюють самим сталим; найменше з $\dfrac14,\dfrac13,\dfrac12,1,\sqrt2$ — це $\dfrac14$. Відповідь: \textbf{А}.}
\zadtask{Укажіть інтеграл, який має \textit{найбільше} значення.

\par\vspace{0.2cm}
\begin{flushleft}
\textbf{А} \quad $\displaystyle\int_0^1 \dfrac{1}{3}\,dx$ \\[0.25cm]
\textbf{Б} \quad $\displaystyle\int_0^1 \dfrac{1}{2}\,dx$ \\[0.25cm]
\textbf{В} \quad $\displaystyle\int_0^1 \sqrt{2}\,dx$ \\[0.25cm]
\textbf{Г} \quad $\displaystyle\int_0^1 1\,dx$ \\[0.25cm]
\textbf{Д} \quad $\displaystyle\int_0^1 \dfrac{2}{5}\,dx$
\end{flushleft}
}
\zadtask{Укажіть інтеграл, який має \textit{найбільше} значення.

\par\vspace{0.2cm}
\begin{flushleft}
\textbf{А} \quad $\displaystyle\int_0^2 0{,}3\,dx$ \\[0.25cm]
\textbf{Б} \quad $\displaystyle\int_0^2 \dfrac{1}{4}\,dx$ \\[0.25cm]
\textbf{В} \quad $\displaystyle\int_0^2 \dfrac{2}{5}\,dx$ \\[0.25cm]
\textbf{Г} \quad $\displaystyle\int_0^2 \dfrac{1}{2}\,dx$ \\[0.25cm]
\textbf{Д} \quad $\displaystyle\int_0^2 1\,dx$
\end{flushleft}
}
\taskBlock{Арифметична прогресія --- сесія 1.06, завдання №5}
\zadtask{В арифметичній прогресії один із членів дорівнює $27$, а різниця прогресії дорівнює $5$. Яке з наведених чисел може бути членом цієї прогресії?}
\answerTable{$49$}{$50$}{$51$}{$52$}{$53$}
\solution{Члени прогресії відрізняються від $27$ на кратне $5$, тобто дають остачу $2$ при діленні на $5$. Серед варіантів лише $52=27+5\cdot 5$. Відповідь: \textbf{Г}.}
\zadtask{В арифметичній прогресії один із членів дорівнює $23$, а різниця прогресії дорівнює $4$. Яке з наведених чисел може бути членом цієї прогресії?}
\answerTable{$45$}{$46$}{$47$}{$48$}{$49$}
\zadtask{У геометричній прогресії один із членів дорівнює $16$, а знаменник прогресії дорівнює $2$. Яке з наведених чисел може бути членом цієї прогресії?}
\answerTable{$24$}{$48$}{$64$}{$96$}{$100$}
\taskBlock{Розпізнавання графіка квадратичної функції --- сесія 1.06, завдання №8}
\zadtask{Графік якої з наведених функцій зображено на рисунку?

\vspace{0.2cm}
\noindent
\begin{minipage}[c]{0.58\textwidth}
\noindent
\textbf{А} \quad $y = (x+1)^2$\\[0.1cm]
\textbf{Б} \quad $y = -(x-1)^2$\\[0.1cm]
\textbf{В} \quad $y = -x^2 + 1$\\[0.1cm]
\textbf{Г} \quad $y = x^2 - 1$\\[0.1cm]
\textbf{Д} \quad $y = -x^2 - 1$
\end{minipage}%
\hfill
\begin{minipage}[c]{0.40\textwidth}
\centering
\begin{tikzpicture}[scale=0.7, thick]
    \draw[->] (-2.3,0) -- (2.3,0) node[right] {\small $x$};
    \draw[->] (0,-2.6) -- (0,1.6) node[above] {\small $y$};
    \draw (1,-0.08) -- (1,0.08) node[above=1pt] {\scriptsize $1$};
    \draw (-0.08,1) -- (0.08,1) node[right=1pt] {\scriptsize $1$};
    \node[below left] at (0,0) {\scriptsize $0$};
    \draw[mainGreen, line width=1.2pt, domain=-1.55:1.55, samples=80] plot (\x, {-(\x)^2 - 1});
\end{tikzpicture}
\end{minipage}
\par\vspace{0.2cm}
}
\solution{Парабола вітками вниз (старший коефіцієнт від'ємний) з вершиною в точці $(0;-1)$. Цьому відповідає $y=-x^2-1$. Відповідь: \textbf{Д}.}
\zadtask{Графік якої з наведених функцій зображено на рисунку?

\vspace{0.2cm}
\noindent
\begin{minipage}[c]{0.58\textwidth}
\noindent
\textbf{А} \quad $y = (x-1)^2$\\[0.1cm]
\textbf{Б} \quad $y = x^2 + 1$\\[0.1cm]
\textbf{В} \quad $y = -x^2 + 1$\\[0.1cm]
\textbf{Г} \quad $y = -(x+1)^2$\\[0.1cm]
\textbf{Д} \quad $y = x^2 - 1$
\end{minipage}%
\hfill
\begin{minipage}[c]{0.40\textwidth}
\centering
\begin{tikzpicture}[scale=0.7, thick]
    \draw[->] (-2.3,0) -- (2.3,0) node[right] {\small $x$};
    \draw[->] (0,-1.6) -- (0,2.0) node[above] {\small $y$};
    \draw (1,-0.08) -- (1,0.08) node[below=1pt] {\scriptsize $1$};
    \draw (-0.08,1) -- (0.08,1) node[left=1pt] {\scriptsize $1$};
    \node[below left] at (0,0) {\scriptsize $0$};
    \draw[mainGreen, line width=1.2pt, domain=-1.55:1.55, samples=80] plot (\x, {-(\x)^2 + 1});
\end{tikzpicture}
\end{minipage}
\par\vspace{0.2cm}
}
\zadtask{Графік якої з наведених функцій зображено на рисунку?

\vspace{0.2cm}
\noindent
\begin{minipage}[c]{0.58\textwidth}
\noindent
\textbf{А} \quad $y = x^2 - 1$\\[0.1cm]
\textbf{Б} \quad $y = (x-1)^2$\\[0.1cm]
\textbf{В} \quad $y = -x^2 + 1$\\[0.1cm]
\textbf{Г} \quad $y = x^2 + 1$\\[0.1cm]
\textbf{Д} \quad $y = -x^2 - 1$
\end{minipage}%
\hfill
\begin{minipage}[c]{0.40\textwidth}
\centering
\begin{tikzpicture}[scale=0.7, thick]
    \draw[->] (-2.3,0) -- (2.3,0) node[right] {\small $x$};
    \draw[->] (0,-1.6) -- (0,2.6) node[above] {\small $y$};
    \draw (1,-0.08) -- (1,0.08) node[below right=0pt] {\scriptsize $1$};
    \draw (-0.08,1) -- (0.08,1) node[left=1pt] {\scriptsize $1$};
    \node[below left] at (0,0) {\scriptsize $0$};
    \draw[mainGreen, line width=1.2pt, domain=-1.8:1.8, samples=80] plot (\x, {(\x)^2 - 1});
\end{tikzpicture}
\end{minipage}
\par\vspace{0.2cm}
}
\taskBlock{Порівняння значень функції --- сесія 2.06, завдання №8}
\zadtask{Задано функцію $f(x) = x^2$. Укажіть правильну нерівність.}
\answerTable{$f(4) < 0$}{$f(-4) < f(4)$}{$f(-1) > f(2)$}{$f(-1) < 0$}{$f(3) < f(4)$}
\solution{Для $f(x)=x^2$ маємо $f(3)=9$, $f(4)=16$, тож $f(3)<f(4)$ -- правильно. Відповідь: \textbf{Д}.}
\zadtask{Задано функцію $f(x) = x^3$. Укажіть правильну нерівність.}
\answerTable{$f(2) < 0$}{$f(-1) > f(1)$}{$f(3) < f(-3)$}{$f(-2) < f(1)$}{$f(0) > f(1)$}
\zadtask{Задано функцію $f(x) = |x|$. Укажіть правильну нерівність.}
\answerTable{$f(-3) < 0$}{$f(-5) > f(2)$}{$f(0) > f(1)$}{$f(4) < f(-4)$}{$f(-2) < f(2)$}
\taskBlock{Сума арифметичної прогресії --- сесія 2.06, завдання №11}
\zadtask{Задано арифметичну прогресію $(a_n)$: $-2;\ -1{,}6;\ -1{,}2;\ \ldots$\ Знайдіть суму перших $21$ членів цієї прогресії.}
\answerTable{$42$}{$-42$}{$21$}{$84$}{$4{,}2$}
\solution{$a_1=-2$, $d=0{,}4$, тож $a_{21}=-2+20\cdot0{,}4=6$. Сума $S_{21}=\dfrac{a_1+a_{21}}{2}\cdot21=\dfrac{-2+6}{2}\cdot21=42$. Відповідь: \textbf{А}.}
\zadtask{Задано арифметичну прогресію $(a_n)$: $-3;\ -2{,}5;\ -2;\ \ldots$\ Знайдіть суму перших $25$ членів цієї прогресії.}
\answerTable{$-75$}{$37{,}5$}{$75$}{$150$}{$7{,}5$}
\zadtask{Задано арифметичну прогресію $(a_n)$: $4;\ 3{,}4;\ 2{,}8;\ \ldots$\ Знайдіть суму перших $11$ членів цієї прогресії.}
\answerTable{$-11$}{$11$}{$22$}{$5{,}5$}{$110$}
\taskBlock{Властивості непарної функції за графіком --- сесія 3.06, завдання №8}
\noindent
\begin{minipage}[c]{0.56\textwidth}
\zadnum У прямокутній системі координат зображено фрагмент графіка \textit{непарної} функції $y = f(x)$ та точки $K$, $L$, $M$, $N$, $P$. Через яку з цих точок проходить графік функції $y = f(x)$?
\end{minipage}%
\hfill
\begin{minipage}[c]{0.40\textwidth}
\centering
\begin{tikzpicture}[scale=0.5, thick]
    \draw[gray!35, very thin] (-3.3,-3.3) grid (3.3,3.3);
    \draw[->] (-3.5,0) -- (3.5,0) node[right, font=\scriptsize] {$x$};
    \draw[->] (0,-3.5) -- (0,3.5) node[above, font=\scriptsize] {$y$};
    \node[below left, font=\tiny] at (0,0) {$0$};
    \node[font=\tiny] at (1,-0.35) {$1$};
    \node[font=\tiny] at (-0.35,1) {$1$};
    % фрагмент у I чверті: відрізок від (1,1) до (2,3) приблизно
    \draw[mainGreen, line width=1.2pt] (1,1) -- (2,3);
    % точки
    \fill (-3,2) circle (2.4pt); \node[left, font=\tiny] at (-3,2) {$L$};
    \fill (-2.6,1.4) circle (2.4pt); \node[left, font=\tiny] at (-2.6,1.4) {$K$};
    \fill (-2,-3) circle (2.4pt); \node[left, font=\tiny] at (-2,-3) {$M$};
    \fill (2,-2.4) circle (2.4pt); \node[right, font=\tiny] at (2,-2.4) {$N$};
    \fill (-3,-3.1) circle (2.4pt); \node[left, font=\tiny] at (-3,-3.1) {$P$};
\end{tikzpicture}
\end{minipage}
\par\vspace{0.2cm}
\answerTable{$K$}{$L$}{$M$}{$N$}{$P$}
\solution{Графік непарної функції симетричний відносно початку координат. Точка, симетрична до точки $(2;\,3)$ цього фрагмента,~--- це $(-2;\,-3)$, тобто точка $M$. Відповідь: \textbf{В} ($M$).}
\noindent
\begin{minipage}[c]{0.56\textwidth}
\zadnum У прямокутній системі координат зображено фрагмент графіка \textit{непарної} функції $y = f(x)$ та точки $K$, $L$, $M$, $N$, $P$. Через яку з цих точок проходить графік функції $y = f(x)$?
\end{minipage}%
\hfill
\begin{minipage}[c]{0.40\textwidth}
\centering
\begin{tikzpicture}[scale=0.5, thick]
    \draw[gray!35, very thin] (-3.3,-3.3) grid (3.3,3.3);
    \draw[->] (-3.5,0) -- (3.5,0) node[right, font=\scriptsize] {$x$};
    \draw[->] (0,-3.5) -- (0,3.5) node[above, font=\scriptsize] {$y$};
    \node[below left, font=\tiny] at (0,0) {$0$};
    \node[font=\tiny] at (1,-0.35) {$1$};
    \node[font=\tiny] at (-0.35,1) {$1$};
    % фрагмент у I чверті: відрізок від (1,3) до (3,1)
    \draw[mainGreen, line width=1.2pt] (1,3) -- (3,1);
    % точки
    \fill (-3,1) circle (2.4pt); \node[left, font=\tiny] at (-3,1) {$K$};
    \fill (-1,-1) circle (2.4pt); \node[below left, font=\tiny] at (-1,-1) {$L$};
    \fill (-2,-2) circle (2.4pt); \node[left, font=\tiny] at (-2,-2) {$M$};
    \fill (2,-2) circle (2.4pt); \node[right, font=\tiny] at (2,-2) {$N$};
    \fill (-3,-3.1) circle (2.4pt); \node[left, font=\tiny] at (-3,-3.1) {$P$};
\end{tikzpicture}
\end{minipage}
\par\vspace{0.2cm}
\answerTable{$K$}{$L$}{$M$}{$N$}{$P$}
\noindent
\begin{minipage}[c]{0.56\textwidth}
\zadnum У прямокутній системі координат зображено фрагмент графіка \textit{парної} функції $y = f(x)$ та точки $K$, $L$, $M$, $N$, $P$. Через яку з цих точок проходить графік функції $y = f(x)$?
\end{minipage}%
\hfill
\begin{minipage}[c]{0.40\textwidth}
\centering
\begin{tikzpicture}[scale=0.5, thick]
    \draw[gray!35, very thin] (-3.3,-3.3) grid (3.3,3.3);
    \draw[->] (-3.5,0) -- (3.5,0) node[right, font=\scriptsize] {$x$};
    \draw[->] (0,-3.5) -- (0,3.5) node[above, font=\scriptsize] {$y$};
    \node[below left, font=\tiny] at (0,0) {$0$};
    \node[font=\tiny] at (1,-0.35) {$1$};
    \node[font=\tiny] at (-0.35,1) {$1$};
    % фрагмент у II чверті: відрізок від (-3,3) до (-1,1)
    \draw[mainGreen, line width=1.2pt] (-3,3) -- (-1,1);
    % точки
    \fill (2,2) circle (2.4pt); \node[right, font=\tiny] at (2,2) {$K$};
    \fill (2,-2) circle (2.4pt); \node[right, font=\tiny] at (2,-2) {$L$};
    \fill (-2,-2) circle (2.4pt); \node[left, font=\tiny] at (-2,-2) {$M$};
    \fill (1,2) circle (2.4pt); \node[above right, font=\tiny] at (1,2) {$N$};
    \fill (3,-3) circle (2.4pt); \node[right, font=\tiny] at (3,-3) {$P$};
\end{tikzpicture}
\end{minipage}
\par\vspace{0.2cm}
\answerTable{$K$}{$L$}{$M$}{$N$}{$P$}
\taskBlock{Первісна функції --- сесія 3.06, завдання №10}
\zadtask{Визначте загальний вигляд первісних для функції $f(x) = 3x^2$.}
\answerTable{$6x^3 + C$}{$6x + C$}{$3x^3 + C$}{$x^2 + C$}{$x^3 + C$}
\solution{Первісна для $f(x) = 3x^2$ має вигляд $F(x) = 3\cdot\dfrac{x^3}{3} + C = x^3 + C$ (перевірка: $(x^3+C)' = 3x^2$). Відповідь: \textbf{Д} ($x^3 + C$).}
\zadtask{Визначте загальний вигляд первісних для функції $f(x) = 4x^3$.}
\answerTable{$12x^2 + C$}{$x^4 + C$}{$4x^4 + C$}{$x^3 + C$}{$12x^4 + C$}
\zadtask{Для функції $f(x) = 6x^2$ визначте первісну $F(x)$, графік якої проходить через точку $(1;\, 5)$.}
\answerTable{$2x^3 + 3$}{$2x^3 + 5$}{$2x^3$}{$x^3 + 4$}{$12x - 7$}
\taskBlock{Точка мінімуму за графіком функції --- сесія 4.06, завдання №8}
\noindent
\begin{minipage}[c]{0.52\textwidth}
\zadnum На рисунку зображено графік функції $y = f(x)$, визначеної на проміжку $[-4;\, 4]$. Укажіть точку мінімуму цієї функції.
\end{minipage}%
\hfill
\begin{minipage}[c]{0.46\textwidth}
\centering
\begin{tikzpicture}[scale=0.55, thick]
    % Сітка
    \draw[gray!40, thin] (-4.5,-4.5) grid (6.5,6.5);

    % Осі
    \draw[->, thick] (-4.5,0) -- (4.8,0) node[below] {$x$};
    \draw[->, thick] (0,-2.5) -- (0,4.8) node[left] {$y$};

    % Підписи осей
    \node[below left] at (0,0) {$0$};
    \node[below] at (1,0) {$1$};
    \node[below] at (-4,0) {$-4$};
    \node[below] at (4,0) {$4$};
    \node[left] at (0,1) {$1$};

    % Графік функції (гладка крива по ключових точках)
    \draw[mainGreen!80!black, very thick] plot[smooth, tension=0.45] coordinates {
        (-4,-2) (-2.3,5) (0,3) (1.5,1) (3,-1) (4,0)
    };

    % Точки на кінцях
    \fill[mainGreen!80!black] (-4,-2) circle (3pt);
    \fill[mainGreen!80!black] (4,0) circle (3pt);

    % Підпис графіка у білому прямокутнику
    \node[fill=white, inner sep=1pt] at (4, 3.8) {$y = f(x)$};
\end{tikzpicture}
\end{minipage}
\par\vspace{0.4cm}
\answerTable{$-2$}{$1$}{$3$}{$2$}{$5$}
\solution{Точка мінімуму~--- абсциса, у якій функція переходить від спадання до зростання (локальна «впадина»). На графіку це відбувається при $x=3$ (там $y=-1$). Відповідь: \textbf{В}.}
\noindent
\begin{minipage}[c]{0.52\textwidth}
\zadnum На рисунку зображено графік функції $y = f(x)$, визначеної на проміжку $[-4;\, 4]$. Укажіть точку максимуму цієї функції.
\end{minipage}%
\hfill
\begin{minipage}[c]{0.46\textwidth}
\centering
\begin{tikzpicture}[scale=0.55, thick]
    % Сітка
    \draw[gray!40, thin] (-4.5,-3.5) grid (6.5,4.5);

    % Осі
    \draw[->, thick] (-4.5,0) -- (4.8,0) node[below] {$x$};
    \draw[->, thick] (0,-3.0) -- (0,4.3) node[left] {$y$};

    % Підписи осей
    \node[below left] at (0,0) {$0$};
    \node[below] at (1,0) {$1$};
    \node[below] at (-4,0) {$-4$};
    \node[below] at (4,0) {$4$};
    \node[left] at (0,1) {$1$};

    % Графік функції
    \draw[mainGreen!80!black, very thick] plot[smooth, tension=0.45] coordinates {
        (-4,1) (-2,-2) (0,1) (1,3) (2.5,1) (4,-1)
    };

    % Точки на кінцях
    \fill[mainGreen!80!black] (-4,1) circle (3pt);
    \fill[mainGreen!80!black] (4,-1) circle (3pt);

    % Підпис графіка у білому прямокутнику
    \node[fill=white, inner sep=1pt] at (4, 3.6) {$y = f(x)$};
\end{tikzpicture}
\end{minipage}
\par\vspace{0.4cm}
\answerTable{$-2$}{$1$}{$3$}{$-1$}{$4$}
\noindent
\begin{minipage}[c]{0.52\textwidth}
\zadnum На рисунку зображено графік функції $y = f(x)$, визначеної на проміжку $[-4;\, 4]$. Укажіть \textit{найбільше значення} цієї функції.
\end{minipage}%
\hfill
\begin{minipage}[c]{0.46\textwidth}
\centering
\begin{tikzpicture}[scale=0.55, thick]
    % Сітка
    \draw[gray!40, thin] (-4.5,-3.5) grid (6.5,4.5);

    % Осі
    \draw[->, thick] (-4.5,0) -- (4.8,0) node[below] {$x$};
    \draw[->, thick] (0,-3.0) -- (0,4.3) node[left] {$y$};

    % Підписи осей
    \node[below left] at (0,0) {$0$};
    \node[below] at (1,0) {$1$};
    \node[below] at (-4,0) {$-4$};
    \node[below] at (4,0) {$4$};
    \node[left] at (0,1) {$1$};

    % Графік функції
    \draw[mainGreen!80!black, very thick] plot[smooth, tension=0.45] coordinates {
        (-4,0) (-2.5,-2) (-1,1) (1,4) (2.5,2) (4,-1)
    };

    % Точки на кінцях
    \fill[mainGreen!80!black] (-4,0) circle (3pt);
    \fill[mainGreen!80!black] (4,-1) circle (3pt);

    % Підпис графіка у білому прямокутнику
    \node[fill=white, inner sep=1pt] at (4, 3.6) {$y = f(x)$};
\end{tikzpicture}
\end{minipage}
\par\vspace{0.4cm}
\answerTable{$-2$}{$1$}{$2$}{$4$}{$3$}
\taskBlock{Геометрична прогресія --- сесія 4.06, завдання №10}
\zadtask{У геометричній прогресії $(b_n)$ з додатними членами відомо, що $\dfrac{b_3}{b_1} = 25$. Обчисліть $\dfrac{b_4}{b_1}$.}
\answerTable{$5$}{$125$}{$525$}{$300$}{$625$}
\solution{$\dfrac{b_3}{b_1}=q^2=25$, а оскільки члени додатні, $q=5$. Тоді $\dfrac{b_4}{b_1}=q^3=125$. Відповідь: \textbf{Б}.}
\zadtask{У геометричній прогресії $(b_n)$ з додатними членами відомо, що $\dfrac{b_3}{b_1} = 9$. Обчисліть $\dfrac{b_4}{b_1}$.}
\answerTable{$3$}{$12$}{$81$}{$27$}{$729$}
\zadtask{У геометричній прогресії $(b_n)$ з додатними членами відомо, що $\dfrac{b_4}{b_2} = 16$. Обчисліть $\dfrac{b_5}{b_1}$.}
\answerTable{$4$}{$64$}{$256$}{$48$}{$32$}
\taskBlock{Графік функції та точка екстремуму зсунутої функції --- сесія 5.06, завдання №8}
\noindent
\begin{minipage}[c]{0.56\textwidth}
\zadnum На рисунку зображено графік функції $y = f(x)$, визначеної на проміжку $[-5;\,5]$. Визначте точку екстремуму функції $y = f(x - 3)$.
\end{minipage}%
\hfill
\begin{minipage}[c]{0.40\textwidth}
\centering
\begin{tikzpicture}[scale=0.45, thick]
    \draw[gray!30, very thin] (-5.3,-2.3) grid (5.3,4.3);
    \draw[->] (-5.5,0) -- (5.6,0) node[right, font=\scriptsize] {$x$};
    \draw[->] (0,-2.5) -- (0,4.5) node[above, font=\scriptsize] {$y$};
    \node[below left, font=\tiny] at (0,0) {$0$};
    \node[font=\tiny] at (1,-0.4) {$1$}; \node[font=\tiny] at (-5,-0.4) {$-5$}; \node[font=\tiny] at (5,-0.4) {$5$};
    \draw[mainGreen, line width=1.2pt]
        (-5,-2) .. controls (-3,0) and (0,3.5) .. (2,4)
        .. controls (3.5,4) and (4.5,2) .. (5,1.5);
    \node[font=\tiny, mainGreen] at (3.5,4) {$y=f(x)$};
\end{tikzpicture}
\end{minipage}
\par\vspace{0.2cm}
\answerTable{$5$}{$1$}{$-1$}{$2$}{$-5$}
\solution{Графік $y=f(x)$ має екстремум (максимум) у точці $x=2$. Заміна $x\to x-3$ зсуває графік праворуч на $3$, тож точка екстремуму $x=2+3=5$. Відповідь: \textbf{А}.}
\noindent
\begin{minipage}[c]{0.56\textwidth}
\zadnum На рисунку зображено графік функції $y = f(x)$, визначеної на проміжку $[-5;\,5]$. Визначте точку екстремуму функції $y = f(x - 2)$.
\end{minipage}%
\hfill
\begin{minipage}[c]{0.40\textwidth}
\centering
\begin{tikzpicture}[scale=0.45, thick]
    \draw[gray!30, very thin] (-5.3,-2.3) grid (5.3,4.3);
    \draw[->] (-5.5,0) -- (5.6,0) node[right, font=\scriptsize] {$x$};
    \draw[->] (0,-2.5) -- (0,4.5) node[above, font=\scriptsize] {$y$};
    \node[below left, font=\tiny] at (0,0) {$0$};
    \node[font=\tiny] at (1,-0.4) {$1$}; \node[font=\tiny] at (-5,-0.4) {$-5$}; \node[font=\tiny] at (5,-0.4) {$5$};
    \draw[mainGreen, line width=1.2pt]
        (-5,-2) .. controls (-3,1) and (0,3.6) .. (1,4)
        .. controls (2.5,4) and (4,1.5) .. (5,0.8);
    \node[font=\tiny, mainGreen] at (3.4,3.6) {$y=f(x)$};
\end{tikzpicture}
\end{minipage}
\par\vspace{0.2cm}
\answerTable{$1$}{$2$}{$-1$}{$3$}{$-3$}
\noindent
\begin{minipage}[c]{0.56\textwidth}
\zadnum На рисунку зображено графік функції $y = f(x)$, визначеної на проміжку $[-5;\,5]$. Визначте точку екстремуму функції $y = f(x) + 3$.
\end{minipage}%
\hfill
\begin{minipage}[c]{0.40\textwidth}
\centering
\begin{tikzpicture}[scale=0.45, thick]
    \draw[gray!30, very thin] (-5.3,-2.3) grid (5.3,4.3);
    \draw[->] (-5.5,0) -- (5.6,0) node[right, font=\scriptsize] {$x$};
    \draw[->] (0,-2.5) -- (0,4.5) node[above, font=\scriptsize] {$y$};
    \node[below left, font=\tiny] at (0,0) {$0$};
    \node[font=\tiny] at (1,-0.4) {$1$}; \node[font=\tiny] at (-5,-0.4) {$-5$}; \node[font=\tiny] at (5,-0.4) {$5$};
    \draw[mainGreen, line width=1.2pt]
        (-5,3.5) .. controls (-3.5,0) and (-2,-2) .. (-1,-2)
        .. controls (1,-2) and (3.5,1.5) .. (5,4);
    \node[font=\tiny, mainGreen] at (-3.6,3.2) {$y=f(x)$};
\end{tikzpicture}
\end{minipage}
\par\vspace{0.2cm}
\answerTable{$-1$}{$2$}{$-4$}{$3$}{$-3$}
\taskBlock{Арифметична прогресія --- сесія 5.06, завдання №15}
\zadtask{В арифметичній прогресії $(a_n)$ відомо, що $a_1 + a_2 = 11{,}2$, а $a_1 + a_4 = 8{,}8$. Визначте різницю $d$ цієї прогресії.}
\answerTable{$-1{,}2$}{$1{,}2$}{$-2{,}4$}{$2{,}4$}{$-0{,}6$}
\solution{$a_1+a_2=2a_1+d=11{,}2$, $a_1+a_4=2a_1+3d=8{,}8$. Віднявши перше від другого: $2d=-2{,}4$, тож $d=-1{,}2$. Відповідь: \textbf{А}.}
\zadtask{В арифметичній прогресії $(a_n)$ відомо, що $a_1 + a_2 = 7{,}4$, а $a_1 + a_4 = 12{,}2$. Визначте різницю $d$ цієї прогресії.}
\answerTable{$-2{,}4$}{$2{,}4$}{$-4{,}8$}{$4{,}8$}{$1{,}2$}
\zadtask{В арифметичній прогресії $(a_n)$ відомо, що $a_3 = 10$, а $a_7 = 2$. Визначте \textit{перший член} $a_1$ цієї прогресії.}
\answerTable{$14$}{$12$}{$6$}{$-14$}{$16$}
\typeTitle{Завдання на встановлення відповідності}
\taskBlock{Відповідність: характеристики функцій та проміжки --- сесія 23.05, завдання №18}
\zadtask{Установіть відповідність між характеристикою функції (1--3) та проміжком (А--Д), якому вона належить. \nmtyear{2026}

\vspace{0.3cm}

\noindent
\begin{minipage}[t]{0.52\textwidth}
\textit{Характеристика} \par \vspace{0.2cm}
\textbf{1} \ \ область визначення функції $y = \sqrt{x - 2}$\\[0.25cm]
\textbf{2} \ \ точка екстремуму функції $y = x^2 + 2x + 3$\\[0.25cm]
\textbf{3} \ \ абсциса точки перетину прямої $y = 6$ з графіком функції $y = x^3$
\end{minipage}
\begin{minipage}[t]{0.24\textwidth}
\textit{Проміжок} \par \vspace{0.2cm}
\textbf{А} \ \ $(-\infty;\, -2)$\\[0.15cm]
\textbf{Б} \ \ $[-2;\, 0)$\\[0.15cm]
\textbf{В} \ \ $[0;\, 1)$\\[0.15cm]
\textbf{Г} \ \ $(1;\, 2)$\\[0.15cm]
\textbf{Д} \ \ $[2;\, +\infty)$
\end{minipage}
\hfill
\begin{minipage}[t]{0.22\textwidth}
\vspace{-0.2cm}
\begin{flushright}\matchingGrid\end{flushright}
\end{minipage}
}
\solution{$1)$ ОДЗ $y=\sqrt{x-2}$: $x\ge 2$, тобто $[2;+\infty)$~(Д). $2)$ Вершина параболи $y=x^2+2x+3$: $x=-\dfrac{2}{2}=-1\in[-2;0)$~(Б). $3)$ $x^3=6 \Rightarrow x=\sqrt[3]{6}\approx 1{,}8\in(1;2)$~(Г). Відповідь: \textbf{1--Д, 2--Б, 3--Г}.}
\zadtask{Установіть відповідність між характеристикою функції (1--3) та проміжком (А--Д), якому вона належить.

\vspace{0.3cm}

\noindent
\begin{minipage}[t]{0.52\textwidth}
\textit{Характеристика} \par \vspace{0.2cm}
\textbf{1} \ \ область визначення функції $y = \sqrt{x - 3}$\\[0.25cm]
\textbf{2} \ \ точка екстремуму функції $y = x^2 + 6x$\\[0.25cm]
\textbf{3} \ \ абсциса точки перетину прямої $y = 15$ з графіком функції $y = x^3$
\end{minipage}
\begin{minipage}[t]{0.24\textwidth}
\textit{Проміжок} \par \vspace{0.2cm}
\textbf{А} \ \ $(-\infty;\, -3)$\\[0.15cm]
\textbf{Б} \ \ $[-3;\, 0)$\\[0.15cm]
\textbf{В} \ \ $[0;\, 2)$\\[0.15cm]
\textbf{Г} \ \ $(2;\, 3)$\\[0.15cm]
\textbf{Д} \ \ $[3;\, +\infty)$
\end{minipage}
\hfill
\begin{minipage}[t]{0.22\textwidth}
\vspace{-0.2cm}
\begin{flushright}\matchingGrid\end{flushright}
\end{minipage}
}
\zadtask{Установіть відповідність між характеристикою функції (1--3) та проміжком (А--Д), якому вона належить.

\vspace{0.3cm}

\noindent
\begin{minipage}[t]{0.52\textwidth}
\textit{Характеристика} \par \vspace{0.2cm}
\textbf{1} \ \ корінь рівняння $2^x = \dfrac{1}{8}$\\[0.25cm]
\textbf{2} \ \ точка мінімуму функції $y = x^2 - x$\\[0.25cm]
\textbf{3} \ \ область визначення функції $y = \sqrt{2x - 8}$
\end{minipage}
\begin{minipage}[t]{0.24\textwidth}
\textit{Проміжок} \par \vspace{0.2cm}
\textbf{А} \ \ $(-\infty;\, -2)$\\[0.15cm]
\textbf{Б} \ \ $[-2;\, 0)$\\[0.15cm]
\textbf{В} \ \ $[0;\, 1)$\\[0.15cm]
\textbf{Г} \ \ $(1;\, 4)$\\[0.15cm]
\textbf{Д} \ \ $[4;\, +\infty)$
\end{minipage}
\hfill
\begin{minipage}[t]{0.22\textwidth}
\vspace{-0.2cm}
\begin{flushright}\matchingGrid\end{flushright}
\end{minipage}
}
\taskBlock{Характеристики функцій (відповідність) --- сесія 30.05, завдання №17}
\zadtask{Установіть відповідність між характеристикою функції (1--3) та проміжком (А--Д), яким є ця характеристика.

\vspace{0.2cm}
\noindent
\begin{minipage}[t]{0.45\textwidth}
\textit{Характеристика функції}\par\vspace{0.15cm}
\textbf{1} \quad область значень функції\\ $y = \sqrt{x+2}$\\[0.25cm]
\textbf{2} \quad область визначення функції \\$y = \lg(x-2)$\\[0.25cm]
\textbf{3} \quad проміжок зростання функції\\ $y = -x^2 + 4x - 7$
\end{minipage}
\hfill
\begin{minipage}[t]{0.3\textwidth}
\textit{Проміжок}\par\vspace{0.15cm}
\textbf{А} \quad $(-\infty;\, +\infty)$\\[0.15cm]
\textbf{Б} \quad $(2;\, +\infty)$\\[0.15cm]
\textbf{В} \quad $[0;\, +\infty)$\\[0.15cm]
\textbf{Г} \quad $(-\infty;\, 2]$\\[0.15cm]
\textbf{Д} \quad $(-\infty;\, 2)$
\end{minipage}
\hfill
\begin{minipage}[t]{0.2\textwidth}
\vspace{0.4cm}
\begin{flushright}\matchingGrid\end{flushright}
\end{minipage}
}
\solution{\textbf{1}: область значень $y=\sqrt{x+2}$ — це $[0;+\infty)$ (В). \textbf{2}: ОДЗ $\lg(x-2)$: $x-2>0$, тобто $(2;+\infty)$ (Б). \textbf{3}: $y=-x^2+4x-7$ — парабола вітками вниз, вершина $x=2$, зростає на $(-\infty;2]$ (Г). Відповідь: \textbf{1~--~В, 2~--~Б, 3~--~Г}.}
\zadtask{Установіть відповідність між характеристикою функції (1--3) та проміжком (А--Д), яким є ця характеристика.

\vspace{0.2cm}
\noindent
\begin{minipage}[t]{0.45\textwidth}
\textit{Характеристика функції}\par\vspace{0.15cm}
\textbf{1} \quad область значень функції\\ $y = \sqrt{x-1}$\\[0.25cm]
\textbf{2} \quad область визначення функції \\$y = \ln(x+3)$\\[0.25cm]
\textbf{3} \quad проміжок зростання функції\\ $y = -x^2 - 6x + 1$
\end{minipage}
\hfill
\begin{minipage}[t]{0.3\textwidth}
\textit{Проміжок}\par\vspace{0.15cm}
\textbf{А} \quad $[0;\, +\infty)$\\[0.15cm]
\textbf{Б} \quad $(-\infty;\, -3]$\\[0.15cm]
\textbf{В} \quad $(-3;\, +\infty)$\\[0.15cm]
\textbf{Г} \quad $(-\infty;\, +\infty)$\\[0.15cm]
\textbf{Д} \quad $[-3;\, +\infty)$
\end{minipage}
\hfill
\begin{minipage}[t]{0.2\textwidth}
\vspace{0.4cm}
\begin{flushright}\matchingGrid\end{flushright}
\end{minipage}
}
\zadtask{Установіть відповідність між характеристикою функції (1--3) та проміжком (А--Д), яким є ця характеристика.

\vspace{0.2cm}
\noindent
\begin{minipage}[t]{0.45\textwidth}
\textit{Характеристика функції}\par\vspace{0.15cm}
\textbf{1} \quad область визначення функції\\ $y = \sqrt{x-5}$\\[0.25cm]
\textbf{2} \quad область визначення функції \\$y = \lg(4-x)$\\[0.25cm]
\textbf{3} \quad проміжок спадання функції\\ $y = -x^2 + 8x - 3$
\end{minipage}
\hfill
\begin{minipage}[t]{0.3\textwidth}
\textit{Проміжок}\par\vspace{0.15cm}
\textbf{А} \quad $[5;\, +\infty)$\\[0.15cm]
\textbf{Б} \quad $(-\infty;\, 4)$\\[0.15cm]
\textbf{В} \quad $(-\infty;\, 4]$\\[0.15cm]
\textbf{Г} \quad $(4;\, +\infty)$\\[0.15cm]
\textbf{Д} \quad $(-\infty;\, +\infty)$
\end{minipage}
\hfill
\begin{minipage}[t]{0.2\textwidth}
\vspace{0.4cm}
\begin{flushright}\matchingGrid\end{flushright}
\end{minipage}
}
\taskBlock{Відповідність: властивості графіків функцій --- сесія 1.06, завдання №17}
\zadtask{Увідповідніть початок речення (1--3) із його закінченням (А--Д) так, щоб утворилося правильне твердження.

\vspace{0.2cm}
\noindent
\begin{minipage}[t]{0.42\textwidth}
\textit{Початок речення}\par\vspace{0.2cm}
\textbf{1} \quad Графік функції $y = \dfrac{2}{x} + 3$\\[0.4cm]
\textbf{2} \quad Графік функції $y = 4^x$\\[0.3cm]
\textbf{3} \quad Графік функції $y = \mathrm{tg}\,x$
\end{minipage}
\hfill
\begin{minipage}[t]{0.42\textwidth}
\textit{Закінчення речення}\par\vspace{0.2cm}
\textbf{А} \quad симетричний відносно початку координат.\\[0.15cm]
\textbf{Б} \quad симетричний відносно осі $y$.\\[0.15cm]
\textbf{В} \quad перетинає лише вісь $x$.\\[0.15cm]
\textbf{Г} \quad перетинає лише вісь $y$.\\[0.15cm]
\textbf{Д} \quad перетинає графік функції $y = x$ лише в одній точці.
\end{minipage}
\hfill
\begin{minipage}[t]{0.14\textwidth}
\vspace{0.4cm}
\begin{flushright}\matchingGrid\end{flushright}
\end{minipage}
}
\solution{\textbf{1)} $y=\dfrac{2}{x}+3$ ($x\ne0$): осі $y$ не перетинає, а вісь $x$ перетинає при $x=-\dfrac{2}{3}$ --- лише вісь $x$ (В); \textbf{2)} $y=4^x>0$: перетинає лише вісь $y$ у точці $(0;1)$ (Г); \textbf{3)} $y=\mathrm{tg}\,x$ --- непарна, симетрична відносно початку координат (А). Відповідь: \textbf{1\,--\,В, 2\,--\,Г, 3\,--\,А}.}
\zadtask{Увідповідніть початок речення (1--3) із його закінченням (А--Д) так, щоб утворилося правильне твердження.

\vspace{0.2cm}
\noindent
\begin{minipage}[t]{0.42\textwidth}
\textit{Початок речення}\par\vspace{0.2cm}
\textbf{1} \quad Графік функції $y = \dfrac{3}{x} - 2$\\[0.4cm]
\textbf{2} \quad Графік функції $y = \left(\dfrac{1}{5}\right)^{x}$\\[0.3cm]
\textbf{3} \quad Графік функції $y = x^3$
\end{minipage}
\hfill
\begin{minipage}[t]{0.42\textwidth}
\textit{Закінчення речення}\par\vspace{0.2cm}
\textbf{А} \quad симетричний відносно початку координат.\\[0.15cm]
\textbf{Б} \quad симетричний відносно осі $y$.\\[0.15cm]
\textbf{В} \quad перетинає лише вісь $x$.\\[0.15cm]
\textbf{Г} \quad перетинає лише вісь $y$.\\[0.15cm]
\textbf{Д} \quad не перетинає жодної з осей координат.
\end{minipage}
\hfill
\begin{minipage}[t]{0.14\textwidth}
\vspace{0.4cm}
\begin{flushright}\matchingGrid\end{flushright}
\end{minipage}
}
\zadtask{Увідповідніть початок речення (1--3) із його закінченням (А--Д) так, щоб утворилося правильне твердження.

\vspace{0.2cm}
\noindent
\begin{minipage}[t]{0.42\textwidth}
\textit{Початок речення}\par\vspace{0.2cm}
\textbf{1} \quad Графік функції $y = -\dfrac{4}{x}$\\[0.4cm]
\textbf{2} \quad Графік функції $y = \log_2 x$\\[0.3cm]
\textbf{3} \quad Графік функції $y = x^2 + 1$
\end{minipage}
\hfill
\begin{minipage}[t]{0.42\textwidth}
\textit{Закінчення речення}\par\vspace{0.2cm}
\textbf{А} \quad симетричний відносно початку координат.\\[0.15cm]
\textbf{Б} \quad симетричний відносно осі $y$.\\[0.15cm]
\textbf{В} \quad перетинає лише вісь $x$.\\[0.15cm]
\textbf{Г} \quad перетинає обидві осі координат.\\[0.15cm]
\textbf{Д} \quad перетинає лише вісь $x$ і вісь $y$ одночасно в початку координат.
\end{minipage}
\hfill
\begin{minipage}[t]{0.14\textwidth}
\vspace{0.4cm}
\begin{flushright}\matchingGrid\end{flushright}
\end{minipage}
}
\taskBlock{Відповідність: функція та координатні чверті --- сесія 2.06, завдання №17}
\zadtask{Установіть відповідність між функцією (1--3) та координатними чвертями (А--Д), у яких розташований її графік.

\vspace{0.3cm}
\noindent
\begin{minipage}[t]{0.4\textwidth}
\textit{Функція}\par\vspace{0.2cm}
\textbf{1} \quad $y = -x - 1$\\[0.4cm]
\textbf{2} \quad $y = \dfrac{1}{x+2}$\\[0.5cm]
\textbf{3} \quad $y = \log_3 x$
\end{minipage}
\hfill
\begin{minipage}[t]{0.32\textwidth}
\textit{Чверті}\par\vspace{0.2cm}
\textbf{А} \quad I і IV\\[0.15cm]
\textbf{Б} \quad I, II і III\\[0.15cm]
\textbf{В} \quad II, III і IV\\[0.15cm]
\textbf{Г} \quad I і III\\[0.15cm]
\textbf{Д} \quad II і IV
\end{minipage}
\hfill
\begin{minipage}[t]{0.2\textwidth}
\vspace{0.5cm}
\begin{flushright}\matchingGrid\end{flushright}
\end{minipage}
}
\solution{$y=-x-1$ -- спадна пряма з від'ємним зсувом, проходить через II, III і IV чверті -- \textbf{В}; $y=\dfrac{1}{x+2}$ -- гіпербола, зміщена ліворуч, її гілки лежать у I, II і III чвертях -- \textbf{Б}; $y=\log_3 x$ зростає при $x>0$, графік у I і IV чвертях -- \textbf{А}. Відповідь: \textbf{1--В, 2--Б, 3--А}.}
\zadtask{Установіть відповідність між функцією (1--3) та координатними чвертями (А--Д), у яких розташований її графік.

\vspace{0.3cm}
\noindent
\begin{minipage}[t]{0.4\textwidth}
\textit{Функція}\par\vspace{0.2cm}
\textbf{1} \quad $y = x + 2$\\[0.4cm]
\textbf{2} \quad $y = -\dfrac{3}{x}$\\[0.5cm]
\textbf{3} \quad $y = \log_{\frac{1}{3}} x$
\end{minipage}
\hfill
\begin{minipage}[t]{0.32\textwidth}
\textit{Чверті}\par\vspace{0.2cm}
\textbf{А} \quad I і IV\\[0.15cm]
\textbf{Б} \quad I, II і III\\[0.15cm]
\textbf{В} \quad II, III і IV\\[0.15cm]
\textbf{Г} \quad I і III\\[0.15cm]
\textbf{Д} \quad II і IV
\end{minipage}
\hfill
\begin{minipage}[t]{0.2\textwidth}
\vspace{0.5cm}
\begin{flushright}\matchingGrid\end{flushright}
\end{minipage}
}
\zadtask{Установіть відповідність між функцією (1--3) та координатними чвертями (А--Д), у яких розташований її графік.

\vspace{0.3cm}
\noindent
\begin{minipage}[t]{0.4\textwidth}
\textit{Функція}\par\vspace{0.2cm}
\textbf{1} \quad $y = 3x$\\[0.4cm]
\textbf{2} \quad $y = -\dfrac{1}{x}$\\[0.5cm]
\textbf{3} \quad $y = -x - 5$
\end{minipage}
\hfill
\begin{minipage}[t]{0.32\textwidth}
\textit{Чверті}\par\vspace{0.2cm}
\textbf{А} \quad I і IV\\[0.15cm]
\textbf{Б} \quad I, II і III\\[0.15cm]
\textbf{В} \quad II, III і IV\\[0.15cm]
\textbf{Г} \quad I і III\\[0.15cm]
\textbf{Д} \quad II і IV
\end{minipage}
\hfill
\begin{minipage}[t]{0.2\textwidth}
\vspace{0.5cm}
\begin{flushright}\matchingGrid\end{flushright}
\end{minipage}
}
\taskBlock{Спільні точки графіків (відповідність) --- сесія 3.06, завдання №17}
\noindent\zadnum До кожного початку речення (1--3) доберіть його закінчення (А--Д) так, щоб утворилося правильне твердження.

\vspace{0.3cm}
\noindent
\begin{minipage}[t]{0.42\textwidth}
\textit{Початок речення}\par\vspace{0.2cm}
\textbf{1} \quad Графіки функцій $y = 2x - 3$ та $y = x + 1$\\[0.3cm]
\textbf{2} \quad Графіки функцій $y = x^2$ та $y = x^2 - 2$\\[0.3cm]
\textbf{3} \quad Графік функції $y = \dfrac{2}{x}$ та графік рівняння $x^2 + y^2 = 4$
\end{minipage}
\hfill
\begin{minipage}[t]{0.42\textwidth}
\textit{Закінчення речення}\par\vspace{0.2cm}
\textbf{А} \quad не мають жодної спільної точки.\\[0.15cm]
\textbf{Б} \quad мають єдину спільну точку.\\[0.15cm]
\textbf{В} \quad мають лише дві спільні точки.\\[0.15cm]
\textbf{Г} \quad мають лише три спільні точки.\\[0.15cm]
\textbf{Д} \quad мають чотири спільні точки.
\end{minipage}
\hfill
\begin{minipage}[t]{0.14\textwidth}
\vspace{0.5cm}
\begin{flushright}\matchingGrid\end{flushright}
\end{minipage}

\solution{\textbf{1}: прямі $y=2x-3$ і $y=x+1$ мають різні кутові коефіцієнти~--- одна спільна точка (Б). \textbf{2}: параболи $y=x^2$ і $y=x^2-2$ зсунуті по вертикалі~--- спільних точок немає (А). \textbf{3}: гіпербола $y=\dfrac{2}{x}$ і коло $x^2+y^2=4$ перетинаються у двох точках (В). Відповідь: \textbf{1~--~Б, 2~--~А, 3~--~В}.}

\vspace{0.5cm}
\noindent\zadnum До кожного початку речення (1--3) доберіть його закінчення (А--Д) так, щоб утворилося правильне твердження.

\vspace{0.3cm}
\noindent
\begin{minipage}[t]{0.42\textwidth}
\textit{Початок речення}\par\vspace{0.2cm}
\textbf{1} \quad Графіки функцій $y = 3x + 1$ та $y = 3x - 2$\\[0.3cm]
\textbf{2} \quad Графіки функцій $y = x^2$ та $y = 4$\\[0.3cm]
\textbf{3} \quad Графіки функцій $y = x^2$ та $y = 2x - 1$
\end{minipage}
\hfill
\begin{minipage}[t]{0.42\textwidth}
\textit{Закінчення речення}\par\vspace{0.2cm}
\textbf{А} \quad не мають жодної спільної точки.\\[0.15cm]
\textbf{Б} \quad мають єдину спільну точку.\\[0.15cm]
\textbf{В} \quad мають лише дві спільні точки.\\[0.15cm]
\textbf{Г} \quad мають лише три спільні точки.\\[0.15cm]
\textbf{Д} \quad мають чотири спільні точки.
\end{minipage}
\hfill
\begin{minipage}[t]{0.14\textwidth}
\vspace{0.5cm}
\begin{flushright}\matchingGrid\end{flushright}
\end{minipage}

\vspace{0.5cm}
\noindent\zadnum До кожного початку речення (1--3) доберіть його закінчення (А--Д) так, щоб утворилося правильне твердження.

\vspace{0.3cm}
\noindent
\begin{minipage}[t]{0.42\textwidth}
\textit{Початок речення}\par\vspace{0.2cm}
\textbf{1} \quad Графіки функцій $y = \dfrac{1}{x}$ та $y = x$\\[0.3cm]
\textbf{2} \quad Графіки функцій $y = \sqrt{x}$ та $y = -1$\\[0.3cm]
\textbf{3} \quad Графік рівняння $x^2 + y^2 = 9$ та графік функції $y = 3$
\end{minipage}
\hfill
\begin{minipage}[t]{0.42\textwidth}
\textit{Закінчення речення}\par\vspace{0.2cm}
\textbf{А} \quad не мають жодної спільної точки.\\[0.15cm]
\textbf{Б} \quad мають єдину спільну точку.\\[0.15cm]
\textbf{В} \quad мають лише дві спільні точки.\\[0.15cm]
\textbf{Г} \quad мають лише три спільні точки.\\[0.15cm]
\textbf{Д} \quad мають чотири спільні точки.
\end{minipage}
\hfill
\begin{minipage}[t]{0.14\textwidth}
\vspace{0.5cm}
\begin{flushright}\matchingGrid\end{flushright}
\end{minipage}
\taskBlock{Найбільше значення функції на проміжку (відповідність) --- сесія 4.06, завдання №17}
\zadtask{Установіть відповідність між функцією, заданою на проміжку (1--3), та її найбільшим значенням (А--Д) на цьому проміжку.

\vspace{0.3cm}
\noindent
\begin{minipage}[t]{0.48\textwidth}
\textit{Функція}\par\vspace{0.2cm}
\textbf{1} \quad $y = \sin 4x + 2,\ x \in (-\infty; +\infty)$\\[0.35cm]
\textbf{2} \quad $y = 4 - x^2,\ x \in [0; 2]$\\[0.35cm]
\textbf{3} \quad $y = 4 + \sqrt{x},\ x \in [0; 1]$
\end{minipage}
\hfill
\begin{minipage}[t]{0.32\textwidth}
\textit{Найбільше значення}\par\vspace{0.2cm}
\textbf{А} \quad $0$\\[0.15cm]
\textbf{Б} \quad $3$\\[0.15cm]
\textbf{В} \quad $4$\\[0.15cm]
\textbf{Г} \quad $5$\\[0.15cm]
\textbf{Д} \quad $6$
\end{minipage}
\hfill
\begin{minipage}[t]{0.16\textwidth}
\vspace{0.4cm}
\begin{flushright}\matchingGrid\end{flushright}
\end{minipage}
}
\solution{1: $\sin4x\leqslant1$, тож найбільше $y=1+2=3$ (Б). 2: $y=4-x^2$ спадає на $[0;2]$, найбільше при $x=0$: $y=4$ (В). 3: $y=4+\sqrt{x}$ зростає на $[0;1]$, найбільше при $x=1$: $y=5$ (Г). Відповідь: \textbf{1~--~Б, 2~--~В, 3~--~Г}.}
\zadtask{Установіть відповідність між функцією, заданою на проміжку (1--3), та її найбільшим значенням (А--Д) на цьому проміжку.

\vspace{0.3cm}
\noindent
\begin{minipage}[t]{0.48\textwidth}
\textit{Функція}\par\vspace{0.2cm}
\textbf{1} \quad $y = \cos 3x + 4,\ x \in (-\infty; +\infty)$\\[0.35cm]
\textbf{2} \quad $y = 9 - x^2,\ x \in [1; 3]$\\[0.35cm]
\textbf{3} \quad $y = 1 + \sqrt{x},\ x \in [0; 4]$
\end{minipage}
\hfill
\begin{minipage}[t]{0.32\textwidth}
\textit{Найбільше значення}\par\vspace{0.2cm}
\textbf{А} \quad $3$\\[0.15cm]
\textbf{Б} \quad $5$\\[0.15cm]
\textbf{В} \quad $8$\\[0.15cm]
\textbf{Г} \quad $9$\\[0.15cm]
\textbf{Д} \quad $10$
\end{minipage}
\hfill
\begin{minipage}[t]{0.16\textwidth}
\vspace{0.4cm}
\begin{flushright}\matchingGrid\end{flushright}
\end{minipage}
}
\zadtask{Установіть відповідність між функцією, заданою на проміжку (1--3), та її \textit{найменшим} значенням (А--Д) на цьому проміжку.

\vspace{0.3cm}
\noindent
\begin{minipage}[t]{0.48\textwidth}
\textit{Функція}\par\vspace{0.2cm}
\textbf{1} \quad $y = 2\sin x + 3,\ x \in (-\infty; +\infty)$\\[0.35cm]
\textbf{2} \quad $y = 6 - x^2,\ x \in [0; 1]$\\[0.35cm]
\textbf{3} \quad $y = \sqrt{x} - 1,\ x \in [1; 4]$
\end{minipage}
\hfill
\begin{minipage}[t]{0.32\textwidth}
\textit{Найменше значення}\par\vspace{0.2cm}
\textbf{А} \quad $1$\\[0.15cm]
\textbf{Б} \quad $0$\\[0.15cm]
\textbf{В} \quad $5$\\[0.15cm]
\textbf{Г} \quad $3$\\[0.15cm]
\textbf{Д} \quad $7$
\end{minipage}
\hfill
\begin{minipage}[t]{0.16\textwidth}
\vspace{0.4cm}
\begin{flushright}\matchingGrid\end{flushright}
\end{minipage}
}
\taskBlock{Властивості функцій (відповідність) --- сесія 5.06, завдання №17}
\zadnum Доберіть до функції (1--3) її властивість (А--Д).

\vspace{0.3cm}
\noindent
\begin{minipage}[t]{0.34\textwidth}
\textit{Функція}\par\vspace{0.2cm}
\textbf{1} \quad $y = x^4 - x^2$\\[0.35cm]
\textbf{2} \quad $y = \log_{0{,}2} x$\\[0.35cm]
\textbf{3} \quad $y = -\dfrac{1}{x}$
\end{minipage}
\hfill
\begin{minipage}[t]{0.52\textwidth}
\textit{Властивість функції}\par\vspace{0.2cm}
\textbf{А} \quad область визначення збігається з областю значень\\[0.12cm]
\textbf{Б} \quad є спадною\\[0.12cm]
\textbf{В} \quad має лише три нулі\\[0.12cm]
\textbf{Г} \quad має лише два нулі\\[0.12cm]
\textbf{Д} \quad графік розташований у I та III чвертях
\end{minipage}
\hfill
\begin{minipage}[t]{0.12\textwidth}
\vspace{0.5cm}
\begin{flushright}\matchingGrid\end{flushright}
\end{minipage}
\par\vspace{0.5cm}
\solution{$1$: $y=x^4-x^2=x^2(x^2-1)$ має нулі $x=0;\,\pm1$, тобто три нулі (В); $2$: $y=\log_{0{,}2}x$ при основі $0<0{,}2<1$ спадна (Б); $3$: $y=-\dfrac1x$ має область визначення $x\ne0$, що збігається з областю значень (А). Відповідь: \textbf{1~--~В, 2~--~Б, 3~--~А}.}
\zadnum Доберіть до функції (1--3) її властивість (А--Д).

\vspace{0.3cm}
\noindent
\begin{minipage}[t]{0.34\textwidth}
\textit{Функція}\par\vspace{0.2cm}
\textbf{1} \quad $y = x^4 - 9x^2$\\[0.35cm]
\textbf{2} \quad $y = 0{,}3^x$\\[0.35cm]
\textbf{3} \quad $y = -\dfrac{5}{x}$
\end{minipage}
\hfill
\begin{minipage}[t]{0.52\textwidth}
\textit{Властивість функції}\par\vspace{0.2cm}
\textbf{А} \quad область визначення збігається з областю значень\\[0.12cm]
\textbf{Б} \quad є спадною\\[0.12cm]
\textbf{В} \quad має лише три нулі\\[0.12cm]
\textbf{Г} \quad має лише два нулі\\[0.12cm]
\textbf{Д} \quad графік розташований у I та III чвертях
\end{minipage}
\hfill
\begin{minipage}[t]{0.12\textwidth}
\vspace{0.5cm}
\begin{flushright}\matchingGrid\end{flushright}
\end{minipage}
\par\vspace{0.5cm}
\zadnum Доберіть до функції (1--3) її властивість (А--Д).

\vspace{0.3cm}
\noindent
\begin{minipage}[t]{0.34\textwidth}
\textit{Функція}\par\vspace{0.2cm}
\textbf{1} \quad $y = \cos x$\\[0.35cm]
\textbf{2} \quad $y = x^3$\\[0.35cm]
\textbf{3} \quad $y = x^2 - 4x$
\end{minipage}
\hfill
\begin{minipage}[t]{0.52\textwidth}
\textit{Властивість функції}\par\vspace{0.2cm}
\textbf{А} \quad область визначення збігається з областю значень\\[0.12cm]
\textbf{Б} \quad є спадною\\[0.12cm]
\textbf{В} \quad має лише три нулі\\[0.12cm]
\textbf{Г} \quad має лише два нулі\\[0.12cm]
\textbf{Д} \quad є парною
\end{minipage}
\hfill
\begin{minipage}[t]{0.12\textwidth}
\vspace{0.5cm}
\begin{flushright}\matchingGrid\end{flushright}
\end{minipage}
\par\vspace{0.5cm}
\typeTitle{Завдання з короткою відповіддю}
\taskBlock{Властивості визначеного інтеграла --- сесія 23.05, завдання №20}
\zadtask{Обчисліть $\displaystyle\int_0^3 4\bigl(f(x) + 2\bigr)\,dx$, якщо $\displaystyle\int_0^3 f(x)\,dx = 10$. \nmtyear{2026}}
\nmtAnswerBox
\solution{$\displaystyle\int_0^3 4(f(x)+2)\,dx=4\int_0^3 f(x)\,dx+4\cdot 2\cdot(3-0)=4\cdot 10+24=64$. Відповідь: \textbf{64}.}
\zadtask{Обчисліть $\displaystyle\int_0^5 3\bigl(f(x) - 1\bigr)\,dx$, якщо $\displaystyle\int_0^5 f(x)\,dx = 8$.}
\nmtAnswerBox
\zadtask{Обчисліть $\displaystyle\int_2^7 f(x)\,dx$, якщо $\displaystyle\int_0^7 f(x)\,dx = 15$ та $\displaystyle\int_0^2 f(x)\,dx = 4$.}
\nmtAnswerBox
\taskBlock{Найменше значення функції на відрізку --- сесія 30.05, завдання №19}
\zadtask{Знайдіть \textit{значення $x_0$}, при якому функція $y = 2x^4 - 9x^2 + 5$ набуває \textit{найменшого} значення на проміжку $[-1;\, 5]$.
\nmtAnswerBox}
\solution{$y'=8x^3-18x=2x(4x^2-9)$; на $[-1;5]$ критичні точки $x=0$ і $x=1{,}5$. Значення: $y(-1)=-2$, $y(0)=5$, $y(1{,}5)=-5{,}125$, $y(5)$ велике. Найменше — при $x_0=1{,}5$. Відповідь: $1{,}5$.}
\zadtask{Знайдіть \textit{значення $x_0$}, при якому функція $y = x^4 - 8x^2 + 3$ набуває \textit{найменшого} значення на проміжку $[-1;\, 3]$.
\nmtAnswerBox}
\zadtask{Знайдіть \textit{значення $x_0$}, при якому функція $y = x^4 - 2x^2 + 7$ набуває \textit{найбільшого} значення на проміжку $[0;\, 5]$.
\nmtAnswerBox}
\taskBlock{Обчислення інтеграла --- сесія 1.06, завдання №19}
\zadtask{Обчисліть інтеграл $\displaystyle\int_{\frac{1}{e}}^{\sqrt{e}} \dfrac{dx}{x}$.}
\nmtAnswerBox
\solution{$\displaystyle\int \dfrac{dx}{x}=\ln x$, тож інтеграл $=\ln\sqrt{e}-\ln\dfrac{1}{e}=\dfrac{1}{2}-(-1)=1{,}5$. Відповідь: \textbf{$1{,}5$}.}
\zadtask{Обчисліть інтеграл $\displaystyle\int_{\frac{1}{e^2}}^{e} \dfrac{dx}{x}$.}
\nmtAnswerBox
\zadtask{Обчисліть інтеграл $\displaystyle\int_{0}^{\ln 3} e^{x}\,dx$.}
\nmtAnswerBox
\taskBlock{Визначений інтеграл --- сесія 2.06, завдання №19}
\zadtask{Обчисліть інтеграл $\displaystyle\int_{1}^{2} \left( \dfrac{2}{x^3} + 4x \right) dx$.
\nmtAnswerBox
}
\solution{Первісна: $\displaystyle\int\!\left(2x^{-3}+4x\right)dx=-\dfrac{1}{x^2}+2x^2$. Підставляємо межі: $\left(-\dfrac{1}{4}+8\right)-\left(-1+2\right)=7{,}75-1=6{,}75$.}
\zadtask{Обчисліть інтеграл $\displaystyle\int_{1}^{2} \left( \dfrac{4}{x^3} + 6x \right) dx$.
\nmtAnswerBox
}
\zadtask{Обчисліть інтеграл $\displaystyle\int_{1}^{3} \left( \dfrac{6}{x^2} + 2x \right) dx$.
\nmtAnswerBox
}
\taskBlock{Критичні точки і точка мінімуму --- сесія 3.06, завдання №19}
\zadtask{Визначте критичні точки функції $f(x) = 4x^3 - 21x^2 - 90x + 19$. У відповідь запишіть ту з них, яка є \textit{точкою мінімуму}.}
\nmtAnswerBox
\solution{$f'(x) = 12x^2 - 42x - 90 = 6(2x^2 - 7x - 15)$. Корені рівняння $2x^2-7x-15=0$ через дискримінант: $D = 49 + 120 = 169$, $\sqrt{D}=13$, $x = \dfrac{7\pm13}{4}$, тобто $x_1 = 5$, $x_2 = -1{,}5$. Похідна змінює знак з $-$ на $+$ у точці $x=5$, тож це точка мінімуму. Відповідь: \textbf{$5$}.}
\zadtask{Визначте критичні точки функції $f(x) = 2x^3 - 3x^2 - 36x + 7$. У відповідь запишіть ту з них, яка є \textit{точкою мінімуму}.}
\nmtAnswerBox
\zadtask{Визначте критичні точки функції $f(x) = 2x^3 + 9x^2 - 24x + 1$. У відповідь запишіть ту з них, яка є \textit{точкою максимуму}.}
\nmtAnswerBox
\taskBlock{Визначений інтеграл --- сесія 4.06, завдання №20}
\zadtask{Обчисліть значення інтеграла
\[ \int_0^3 \frac{x^2 + 3x}{x + 3} \,dx. \]
\nmtAnswerBox}
\solution{$\dfrac{x^2+3x}{x+3}=\dfrac{x(x+3)}{x+3}=x$. Тоді $\displaystyle\int_0^3 x\,dx=\dfrac{x^2}{2}\Big|_0^3=\dfrac{9}{2}=4{,}5$. Відповідь: \textbf{$4{,}5$}.}
\zadtask{Обчисліть значення інтеграла
\[ \int_0^4 \frac{x^2 + 5x}{x + 5} \,dx. \]
\nmtAnswerBox}
\zadtask{Обчисліть значення інтеграла
\[ \int_0^3 \frac{x^2 - 9}{x - 3} \,dx. \]
\nmtAnswerBox}
\taskBlock{Значення похідної в точці --- сесія 5.06, завдання №20}
\zadtask{Знайдіть значення похідної функції $f(x) = \dfrac{\sqrt{x}}{3x - 2}$ у точці $x_0 = 4$.}
\nmtAnswerBox
\solution{За правилом частки $f'(x)=\dfrac{\frac{1}{2\sqrt x}(3x-2)-\sqrt x\cdot3}{(3x-2)^2}$. При $x_0=4$: $\sqrt4=2$, $3\cdot4-2=10$; чисельник $=\frac{1}{4}\cdot10-2\cdot3=-3{,}5$, знаменник $=100$. Отже $f'(4)=-0{,}035$. Відповідь: \textbf{$-0{,}035$}.}
\zadtask{Знайдіть значення похідної функції $f(x) = \dfrac{\sqrt{x}}{x + 6}$ у точці $x_0 = 4$.}
\nmtAnswerBox
\zadtask{Знайдіть значення похідної функції $f(x) = (3x - 2)\sqrt{x}$ у точці $x_0 = 4$.}
\nmtAnswerBox
\chapterTitle{ВІДПОВІДІ}
\noindent\small Відповіді наведено у тому ж порядку, що й завдання.\par\vspace{0.3cm}
\ansType{Завдання з вибором однієї відповіді (А--Д)}
\begin{multicols}{3}\noindent
\ansitem{А}
\ansitem{Б}
\ansitem{А}
\ansitem{В}
\ansitem{Б}
\ansitem{А}
\ansitem{Б}
\ansitem{В}
\ansitem{Д}
\ansitem{А}
\ansitem{В}
\ansitem{Д}
\ansitem{Г}
\ansitem{В}
\ansitem{В}
\ansitem{Д}
\ansitem{В}
\ansitem{А}
\ansitem{Д}
\ansitem{Г}
\ansitem{Б}
\ansitem{А}
\ansitem{В}
\ansitem{Б}
\ansitem{В}
\ansitem{В}
\ansitem{А}
\ansitem{Д}
\ansitem{Б}
\ansitem{А}
\ansitem{В}
\ansitem{Б}
\ansitem{Г}
\ansitem{Б}
\ansitem{Г}
\ansitem{В}
\ansitem{А}
\ansitem{Г}
\ansitem{А}
\ansitem{А}
\ansitem{Б}
\ansitem{А}
\end{multicols}
\ansType{Завдання на встановлення відповідності}
\begin{multicols}{3}\noindent
\ansitem{1~--~Д, 2~--~Б, 3~--~Г}
\ansitem{1~--~Д, 2~--~Б, 3~--~Г}
\ansitem{1~--~А, 2~--~В, 3~--~Д}
\ansitem{1~--~В, 2~--~Б, 3~--~Г}
\ansitem{1~--~А, 2~--~В, 3~--~Б}
\ansitem{1~--~А, 2~--~Б, 3~--~Г}
\ansitem{1~--~В, 2~--~Г, 3~--~А}
\ansitem{1~--~В, 2~--~Г, 3~--~А}
\ansitem{1~--~А, 2~--~В, 3~--~Б}
\ansitem{1~--~В, 2~--~Б, 3~--~А}
\ansitem{1~--~Б, 2~--~Д, 3~--~А}
\ansitem{1~--~Г, 2~--~Д, 3~--~В}
\ansitem{1~--~Б, 2~--~А, 3~--~В}
\ansitem{1~--~А, 2~--~В, 3~--~Б}
\ansitem{1~--~В, 2~--~А, 3~--~Б}
\ansitem{1~--~Б, 2~--~В, 3~--~Г}
\ansitem{1~--~Б, 2~--~В, 3~--~А}
\ansitem{1~--~А, 2~--~В, 3~--~Б}
\ansitem{1~--~В, 2~--~Б, 3~--~А}
\ansitem{1~--~В, 2~--~Б, 3~--~А}
\ansitem{1~--~Д, 2~--~А, 3~--~Г}
\end{multicols}
\ansType{Завдання з короткою відповіддю}
\begin{multicols}{3}\noindent
\ansitem{$64$}
\ansitem{$9$}
\ansitem{$11$}
\ansitem{$1{,}5$}
\ansitem{$2$}
\ansitem{$5$}
\ansitem{$1{,}5$}
\ansitem{$3$}
\ansitem{$2$}
\ansitem{$6{,}75$}
\ansitem{$10{,}5$}
\ansitem{$12$}
\ansitem{$5$}
\ansitem{$3$}
\ansitem{$-4$}
\ansitem{$4{,}5$}
\ansitem{$8$}
\ansitem{$13{,}5$}
\ansitem{$-0{,}035$}
\ansitem{$0{,}005$}
\ansitem{$8{,}5$}
\end{multicols}

\end{document}
